\section{From the recovered quotients to the field}
\label{sec:reconstruction}

We have recovered the geometric data needed for the last step.  We now
specify that data and Pop's reconstruction theorems, so that their
application uses exactly the original group input.

For a function field $F/\kappa$, a \emph{divisorial flag} is a sequence
of valuations obtained by taking a prime divisor trivial on $\kappa$,
then a prime divisor on its residue field, and continuing in this way.
The empty flag is included.  The \emph{total decomposition graph}
records, at each such flag, the abelian pro-$\ell$ group of the residue
field, and at each next divisorial valuation its inertia and
decomposition subgroups and the quotient map to the next residue group.
Thus the graph retains the group maps as well as the flags.  At a
one-variable field over algebraically closed constants, its terminal
data are precisely the point inertia lines: the residual groups at
closed points vanish, so point decomposition equals point inertia.

We use the following two results of Pop.
\begin{proposition}[Pop's reconstruction theorems]
\label{reconstruction:pop}
Let $F/\kappa$ and $F'/\kappa'$ be function fields of transcendence
degree greater than one over algebraically closed fields of
characteristic different from $\ell$.
\begin{enumerate}
\item From $\Pi_F^c\to\Pi_F^a$ and the union of its true divisorial
inertia groups, one recovers the total decomposition graph, functorially
under class-two isomorphisms whose abelianizations preserve this union.
\item Suppose an isomorphism between the total decomposition graphs of
$F$ and $F'$ is compatible with all their geometric rational quotient
diagrams.  Then its map on the top abelian groups is induced, up to one
unit in $\Zell$, by an isomorphism of perfect closures in the opposite
direction carrying the constant fields onto each other.  That field
isomorphism is unique modulo Frobenius twists in positive characteristic,
and unique in characteristic zero.
\end{enumerate}
In (2), a geometric rational quotient is restriction to
$\kappa(u)\subset F$ with $F/\kappa(u)$ regular, together with the
point inertia and decomposition data on its target.  Compatibility
means that restriction commutes with the graph isomorphisms and the
induced isomorphisms of these targets.
\end{proposition}
\begin{proof}[Reference and hypotheses]
Part (1) is \cite[Proposition~2.4]{Pop2018}; the surrounding section
assumes transcendence degree greater than one.  Part (2) is the Isom
statement of \cite[Theorem~2.1(2)]{PopBI2012}, restating
\cite[Main Theorem]{PopGraphs2012}.  The stated theorem requires
Bertini-type families of rational quotients; all geometric rational
quotients form such a family.  More explicitly, for algebraically
independent $x,t$ with $x$ separating, the general functions
\[
 ax+t,\qquad \frac{t}{a'x+a},\qquad
 \frac{a''t+a'x+a+1}{t+a'x+a}
\]
in \cite[Fact/Definition~43]{PopGraphs2012} define relatively
algebraically closed rational subfields for the prescribed cofinite
choices of constants.  Every one of these subfields is among the
quotients in (2).  Thus its hypotheses imply precisely the required
Bertini compatibility.  No restriction to finite-field closures or to
dimension greater than two occurs in this Isom statement.
\end{proof}

We emphasize the distinction between the two citations.  The later
\cite[Theorem~1.1]{Pop2018} reconstructs rational quotients from
divisorial inertia under a dimension-greater-than-two hypothesis.
Here those quotients have already been recovered by
Section~\ref{sec:curve-recovery}, so that theorem is unnecessary.

\begin{proof}[Proof of Theorem~\ref{thm:main}]
A field isomorphism $\alpha:K^i\to L^i$ induces a contravariant
isomorphism of absolute Galois groups, well defined up to an inner
automorphism.  Purely inseparable extension does not change these
groups.  Abelianization removes the inner ambiguity, and the induced
class-two map preserves the commutator.  Thus it gives a member of
$\Isom^c(\Pi_L^a,\Pi_K^a)$.  In positive characteristic, absolute
Frobenius on algebraic closures commutes with field automorphisms;
composition with any integer Frobenius power therefore gives the same
allowed orbit.  This proves well-definedness of $\Phi_{K,L}$.

Conversely, let $\varphi:\Pi_L^a\to\Pi_K^a$ be bracket-compatible.
Lemma~\ref{prelim:central-lift} lifts it to a class-two isomorphism.
After choosing Tate identifications, let $\Theta$ be its dual
\eqref{eq:theta}.  Proposition~\ref{local:transport} recovers the common
relative transcendence degree and the quasi-divisorial pairs.
The finite residual tests and the uniform support bound prove the
completed curve-subfield correspondence in
Theorem~\ref{recovery:curves}.  The recognition statements in
Section~\ref{sec:curve-recovery} then show that $\varphi$ preserves true
divisorial inertia and every relatively algebraically closed rational
quotient, including its point lines.

Proposition~\ref{recovery:rational} gives the actual commuting
quotient diagrams in \eqref{eq:recovery-quotient-diagram}.
Their vertical maps are surjective, and their target isomorphisms
$\varphi_E$ match all point inertia lines.  For rational targets these
are the entire point decomposition data.  Thus the same fixed
$\varphi$ is compatible with every geometric rational quotient;
no independent rescaling of the targets is needed.

Part (1) of Proposition~\ref{reconstruction:pop} gives the total
decomposition-graph isomorphism.  The quotient squares also commute on
these graphs.  Indeed, restriction along an actual regular inclusion
$E=k(t)\subset K$ is a geometric rational quotient morphism
\cite[Section~2, Embeddings and Restrictions]{PopBI2012}.
At a flag vertex its map is obtained by restricting to the decomposition
subgroup and passing to the inertia quotient
\cite[Definition/Remark~24 and Definition~25]{PopGraphs2012}.
If the flag restricts trivially to $E$, its inertia is killed, and the
top square descends to the residual square.  Otherwise its nonzero
inertia image lies in a unique point line of $E$, where decomposition
equals inertia and the residue group is trivial.  The same description
applies to $P\subset L$.  Local images can have finite index: saturation
identifies the point line, but the actual subgroup images and maps are
retained.  Applying these restriction and quotient operations successively
along each flag therefore proves
compatibility for the same fixed $\varphi$ and $\varphi_E$.

Part (2), applied to the recovered rational quotients,
gives an isomorphism $\alpha:K^i\to L^i$ with $\alpha(k)=l$ and
$[\varphi]=[\alpha^*]$.  Choices of Tate identifications are absorbed
by the one allowed global unit.  This proves surjectivity.  It also
proves that different characteristics cannot give a bracket-compatible
isomorphism.

Finally, two field isomorphisms with the same Galois-side orbit give
the same compatible graph and quotient data after one global unit
adjustment.  The uniqueness assertion in
Proposition~\ref{reconstruction:pop} identifies them modulo Frobenius
in positive characteristic, and identifies them outright in
characteristic zero.  On perfect closures the Frobenius powers are
exactly the integer powers appearing in the definition of
$\Isom_F^i(K,L)$.  This proves injectivity and completes the theorem.
\end{proof}
